\section{An isomorphism with de Bruijn terms}
\label{sec:de-bruijn}

TODO: fill in the actual de Bruijn development being reused
(source/authorship, what was already in CSLib/Mathlib vs. what was added),
the precise choice of which of the five equivalent alpha
equivalences is quotiented by, and the isomorphism statement and
proof sketch, following Norrish's proof pearl.

De Bruijn indices remove the need for alpha equivalence altogether by
replacing bound variable names with numerals counting binder depth.
Norrish's proof pearl~\cite{norrish2007a} demonstrates that,
despite a reputation for being awkward to reason about, a de Bruijn
representation supports clean, structurally recursive proofs of the
standard metatheory of the lambda calculus once the representation is
set up carefully. We reuse an existing de Bruijn development
and show that it is isomorphic to the quotient of named terms
$\Lam$ by (any one of) the alpha equivalences of
Section~\ref{sec:alpha-equivalences}.

\subsection{De Bruijn terms}
\label{subsec:db-terms}

\begin{definition}[De Bruijn terms]
\label{def:db-terms}
The set $\LamDB$ of de Bruijn lambda terms is generated by
\[
  t \;::=\; n \mid t\, t \mid \lambda.\, t
  \qquad (n \in \mathbb{N}).
\]
\end{definition}

\subsection{The isomorphism}
\label{subsec:isomorphism}

\begin{theorem}
\label{thm:db-isomorphism}
$\Lam / \alphaeq \;\cong\; \LamDB$.
\end{theorem}

